\subsection{Sprint 2}

A Sprint 2, iniciada em 28 de abril, foi marcada pela execução de uma refatoração estrutural ampla da organização de pastas do projeto, tarefa que assumi de forma individual dada minha responsabilidade pela arquitetura do frontend. Esse esforço, que acabou impactando mais de cento e vinte arquivos do repositório, teve como objetivo reorganizar o código segundo os princípios de \textit{feature-based architecture} e de separação entre componentes de apresentação e de lógica definidos desde a Sprint 0, padronizando nomenclaturas e eliminando inconsistências que haviam se acumulado com o crescimento orgânico do número de telas desenvolvidas pela equipe.

Em paralelo a essa refatoração, contribuí com ajustes pontuais, porém relevantes, em telas já existentes: integrei os \textit{insights} de reuniões à tela de detalhe de reunião, adicionei e posteriormente ajustei o campo de sentimento médio (\textit{average feeling}) no detalhe do vendedor conforme o time refinava os requisitos junto aos \textit{stakeholders}, e realizei a limpeza dos seletores de empresa, corrigindo problemas de \textit{layout} que vinham sendo carregados de sprints anteriores. Também exerci meu papel de revisão arquitetural validando e integrando os \textit{pull requests} da tela de detalhe do vendedor e da formatação da taxa de sucesso das reuniões, ambos desenvolvidos por colegas de equipe.

O principal desafio da sprint surgiu justamente como consequência do tamanho da refatoração de pastas: como diversos colegas estavam com \textit{branches} de funcionalidades abertas em paralelo (telas de listagem e detalhe de reuniões, modal de cadastro de vendedores, entre outras), a integração da reestruturação à \textit{branch} de desenvolvimento gerou uma sequência de conflitos de \textit{merge} que precisaram ser resolvidos por cada um dos colegas em suas respectivas tarefas, atrasando a finalização de algumas entregas. Esse episódio me mostrou a importância de comunicar com antecedência mudanças estruturais de grande escala e, quando possível, fragmentá-las em etapas menores para reduzir o impacto sobre o trabalho em andamento do restante da equipe.

Encerrada a Sprint 2, a base de código encontrava-se consideravelmente mais organizada e alinhada à arquitetura planejada, o que facilitou a integração das funcionalidades restantes. Com a estrutura do projeto madura, o foco da equipe se direcionou para a Sprint 3, com a construção das primeiras visualizações de dados do painel administrativo.
